\subsection{Exact ranks of the elliptic and quadratic Prym factors}

\begin{statement}{A.185}{Proved}
{Three explicit Weil-restriction isogenies and a rank contribution of eight}
Use the curves, Prym surfaces, and actual Jacobian isogeny of A.182.
There are actual \(\Q\)-isogenies
\[
\begin{array}{c|c|c|c}
\text{Prym surface}&L& E/L&\operatorname{rank}E(L)\\ \hline
\mathcal P_{-12,1}&\Q(\sqrt{-3})&y^2=x^3-4\sqrt{-3}\,x&1\\
\mathcal P_{18,1}&\Q(\sqrt2)&y^2=x^3+6\sqrt2\,x&2\\
\mathcal P_{3,-2}&\Q(\sqrt{-6})&y^2=x^3+8\sqrt{-6}\,x&2
\end{array}
\qquad
\mathcal P_{A,B}\longrightarrow\operatorname{Res}_{L/\Q}E.
\tag{A.185.1}
\]
The three rational elliptic factors \(E_{-6},E_{18},E_{-12}\) of
A.182 each have rank one.  Consequently the exact rank identity is
\[
 \operatorname{rank}J(C_8)(\Q)
 =8+\sum_{(A,B)\in\mathcal G}
          \operatorname{rank}\operatorname{Jac}(G_{A,B})(\Q)
 \geq14,
\tag{A.185.2}
\]
where \(\mathcal G\) is the six-element set in (A.182.3).
The proof gives all connecting images used in the six elliptic rank
computations.  It supplies no upper bound for the six remaining
genus-two ranks or for the total rank in (A.185.2).
\end{statement}

\begin{proof}
\emph{The quotient maps and their descent.}
Write \(r^2=B/A\), so that the proper plane quartic has affine equation
\(D:w^4=Az(z^2+r^2)\), and put \(L=\Q(r)\).  In the three cases
of the statement this is a quadratic field.  The actual involutions
and their elliptic quotient functions are
\[
\begin{aligned}
 j_\pm(z,w)&=(r^2/z,\ \pm rw/z),\qquad j_+j_-=\iota,\\
 q_\pm:D_L&\longrightarrow E_\pm,
 &X&=Aw^2/z,&Y_\pm&=A^2w(z\pm r)/z,\\
 E_\pm&:\quad Y_\pm^2=X^3\pm2A^3rX.
\end{aligned}
\tag{A.185.3}
\]
Substitution proves the equations and that \(q_\pm\) is fixed by
\(j_\pm\).  For the plus map the inverse field equations are
\[
 z^2-\frac{X^2}{A^3}z+r^2=0,
 \qquad w=\frac{Y_+z}{A^2(z+r)}.
\tag{A.185.4}
\]
Thus the map has degree at most two; the nonidentity involution
\(j_+\) fixes its functions, so the degree is exactly two.
Replacing \(r\) by \(-r\) gives the other map.  Each rational map
extends uniquely to the proper smooth curves.

The pullbacks of the regular elliptic differentials are literally
\[
 q_\pm^*\frac{dX}{2Y_\pm}
       =\frac{(z\mp r)\,dz}{4w^3}.
\tag{A.185.5}
\]
They span the anti-\(\iota\) differential space.  Moreover
\(q_\pm\iota=[-1]q_\pm\), so the divisor pushforwards restrict
to homomorphisms on \(\mathcal P_{A,B,L}\).  The nontrivial
automorphism of \(L/\Q\) exchanges the two maps, with their
displayed target equations.  Adjunction for Weil restriction
therefore gives a \(\Q\)-homomorphism
\[
 \mathcal P_{A,B}\longrightarrow
       \operatorname{Res}_{L/\Q}E_+.
\tag{A.185.6}
\]
Over \(L\) its two components are the displayed pushforwards.
On the anti-invariant cotangent space their matrix, apart from the
common nonzero factor \(1/4\), is
\(\left(\begin{smallmatrix}-r&1\\r&1\end{smallmatrix}\right)\),
of determinant \(-2r\).  Its differential is an isomorphism;
both abelian varieties have dimension two, so (A.185.6) is an
isogeny.  This proves the exact map to which the rank comparison
will be applied.

For each row put \(s^2=d\).  The changes
\((X,Y)=(\lambda^2x,\lambda^3y)\) and the images of the rational
points on \(D\) used in A.182 are as follows:
\[
\begin{array}{c|r|c|c|c|c}
(A,B)&d&r&\lambda&b\text{ in }E_b&\text{image point}\\ \hline
(-12,1)&-3&s/6&2s&-4s&(-2,\,2+2s)\\
(18,1)&2&s/6&3s&6s&(3,\,3+3s)\\
(3,-2)&-6&s/3&s/2&8s&(-2,\,-4+2s)
\end{array}
\tag{A.185.7}
\]
Here \(E_b:y^2=x^3+bx\).  The source points are respectively
\((-1/2,1),(1/3,1),(1,1)\).  In particular every scaling and
point in the table is defined over the stated field.

\emph{The exact two-isogeny calculation.}
For any of the fields and coefficients below, use the literal maps
\[
\begin{aligned}
 \phi_b:E_b&\longrightarrow E_{-4b},
 &(x,y)&\longmapsto(x+b/x,\ y(1-b/x^2)),\\
 \psi_b:E_{-4b}&\longrightarrow E_b,
 &(x,y)&\longmapsto(x/4-b/x,\ y(1+4b/x^2)/8).
\end{aligned}
\tag{A.185.8}
\]
They extend at the omitted points, have kernel \(\{O,(0,0)\}\),
and satisfy \(\psi_b\phi_b=[2]\).  The connecting map for
\(\psi_b\) is
\[
 \alpha_b:E_b(k)\longrightarrow k^\times/k^{\times2},
 \quad O\longmapsto1,\quad(0,0)\longmapsto[b],
 \quad(x,y)\longmapsto[x]\quad(x\ne0).
\tag{A.185.9}
\]
Here is an explicit verification of its kernel and cocycle.
The first coordinate of \(\psi_b(Q)\) is
\((y(Q)/(2x(Q)))^2\).  Conversely, if \(P=(s_0^2,y)\) with
\(s_0\ne0\), then
\[
 x(Q)=2s_0^2+2y/s_0,\qquad y(Q)=2s_0x(Q)
\tag{A.185.10}
\]
satisfies the equation of \(E_{-4b}\) and maps to \(P\).
Changing the square root \(s_0\) changes \(Q\) by the kernel
point.  Thus its actual connecting cocycle is the square-root
cocycle of \(x(P)\).  The exceptional point \((0,0)\) has a
preimage exactly when \(b\) is a square: its two preimages then
have first coordinate \(\pm2\sqrt b\).  This also verifies the
exceptional value in (A.185.9).  The connecting construction proves
that (A.185.9) is a homomorphism, with kernel \(\psi_bE_{-4b}(k)\).

At a prime not dividing \(b\), a nonzero valuation of \(x(P)\)
is even.  If it is positive, \(x^2+b\) is a unit; if negative,
the valuation of \(x(x^2+b)\) is three times that of \(x\).
Both assertions follow from the actual equation of \(E_b\).
The squareclass \(d_0\) occurs only if the proper quartic
\[
       N^2=d_0U^4+(b/d_0)V^4
\tag{A.185.11}
\]
has a point, with the actual map
\(x=d_0(U/V)^2,\ y=d_0UN/V^3\); the endpoint charts give
the exceptional classes already described.  Locally one can
scale to primitive coordinates, with at least one of \(U,V\)
a unit.  When \(v(b)=1\), this gives the following useful exact
necessary test.  If \(v(d_0)=0\), then \(U\) is a unit and
the residue of \(d_0\) must be a square.  If \(v(d_0)=1\),
then \(V\) is a unit and the residue of \(b/d_0\) must be a
square.  In the other chart the right side has odd valuation.

In every calculation below \(E_b(k)[2]\) has order two, and
the nonidentity point of \(E_{-4b}[\psi_b](k)\) has nontrivial
connecting class \([-4b]\) for \(\phi_b\).  The maps (A.185.8)
therefore give the exact index identity
\[
\begin{aligned}
 |E_b(k)/2E_b(k)|
 &=|\alpha_b(E_b(k))|\,
   |E_{-4b}(k)/(\phi_bE_b(k)+E_{-4b}[\psi_b](k))|\\
 &=\frac{|\alpha_b(E_b(k))|\,
              |\alpha_{-4b}(E_{-4b}(k))|}{2},\\
 2^{\operatorname{rank}E_b(k)}
 &=\frac{|\alpha_b(E_b(k))|\,
              |\alpha_{-4b}(E_{-4b}(k))|}{4}.
\end{aligned}
\tag{A.185.12}
\]
The last equality uses finite generation over the stated number
field, exactly \cite[I, Theorem 16.7, p.69]{LegMilneAV2008}.
The condition on two-torsion follows in the three quadratic cases
from the odd valuation of \(-b\) at a prime above three.  It
also holds directly for the three rational coefficients below.

\emph{The field \(\Q(\sqrt{-3})\).}
Put \(s^2=-3\).  Its ring of integers is
\(\Z[(1+s)/2]\).  Nearest integral coefficients in this basis
leave a remainder of norm at most \(3/4<1\), since its norm
form is \(a^2+ab+b^2\).  The ring is norm-Euclidean, its ideal
class group is trivial, and its units are the six roots of unity.
Modulo unit squares they are generated by \(-1\).  The prime
two is inert, and \((s)\) is the prime above three.

The valuation test above consequently confines the connecting
classes for \(b=-4s\) to
\((-1)^\sigma2^es^f\), with \(\sigma,e,f\in\{0,1\}\).
For the dual use the actual scaling
\(E_s\longrightarrow E_{16s}\), \((x,y)\mapsto(4x,8y)\).
It multiplies the first coordinate by a square; the possible
classes for \(E_s\) are just \(\pm1,\pm s\).
At \((s)\), the primitive test following (A.185.11), in the
residue field \(\F_3\), gives
\[
\begin{aligned}
 \alpha_{-4s}(E_{-4s}(k))&\subset\{1,-2,-s,2s\},\\
 \alpha_s(E_s(k))&\subset\{1,s\}.
\end{aligned}
\tag{A.185.13}
\]
For example the tested residue is \((-1)^\sigma2^e\) when
\(f=0\), and \(-4/((-1)^\sigma2^e)\) when \(f=1\).
The point \((-2,2+2s)\) and \((0,0)\) on the first curve
give the independent classes \(-2,-s\); the kernel point on
\(E_s\) gives \(s\).  Thus both inclusions are equalities.
Their orders are four and two, giving rank one by (A.185.12).

\emph{The field \(\Q(\sqrt2)\).}
Now \(s^2=2\) and \(\epsilon=1+s\).  Nearest integral
coefficients in \(\Z[s]\) leave a remainder of absolute norm
at most \(1/2\), so again the class group is trivial.  Its
units are \(\pm\epsilon^n\).  For completeness,
\(\epsilon\) is the least unit with positive real value greater
than one: for \(1<u<\epsilon\), the integer coefficient
\((u+u^{-1})/2\) in the norm-one case lies strictly between
one and two, and the coefficient \((u-u^{-1})/2\) in the
norm-minus-one case lies strictly between zero and one.
Dividing any positive unit by a power of \(\epsilon\) proves
the assertion.  The support primes are \((s)\) and the inert
prime \((3)\); all candidate classes therefore have the form
\((-1)^\sigma\epsilon^es^\delta3^\gamma\), with four binary
exponents.

At three, \(v_3(\pm6s)=1\).  A unit of \(\F_9\) is a
square exactly when its norm to \(\F_3\) is one.
The norms of \(-1,s,\epsilon\) in this residue field are
respectively \(1,1,-1\).  In both branches of the primitive
test this forces \(e=0\), whether \(\gamma=0\) or one.
For \(b=6s\), the positive real embedding has \(b>0\),
so \(x(P)\) is positive and \(\sigma=0\).  The actual
scaled dual is \(E_{-6s}\), via
\((x,y)\mapsto(s^2x,s^3y)\) to \(E_{-24s}\).
At its real embedding with \(s<0\), positivity instead
forces \(\sigma=\delta\).  It follows that
\[
\begin{aligned}
 \alpha_{6s}(E_{6s}(k))&\subset\{1,3,s,3s\},\\
 \alpha_{-6s}(E_{-6s}(k))&\subset\{1,3,-s,-3s\}.
\end{aligned}
\tag{A.185.14}
\]
The points \((3,3+3s)\) and \((3,3-3s)\), respectively,
give the class three.  Their kernel points give \(3s\) and
\(-3s\), since \(2=s^2\).  The valuations at three and
\((s)\) show independence.  Both images have order four,
so this rank is two.  No dyadic sufficiency assertion is used.

\emph{The field \(\Q(\sqrt{-6})\).}
Here \(s^2=-6\), \(\mathcal O_k=\Z[s]\), and the
discriminant is \(-24\).  The exact Minkowski bound
\((2/\pi)\sqrt{24}<4\), from
\cite[Theorem 4.3, p.70]{AgentMilneANT}, shows that every ideal
class is represented by an ideal of norm one, two, or three.
The unique primes of norms two and three satisfy
\[
 \mathfrak p_2^2=(2),\qquad \mathfrak p_3^2=(3),
 \qquad\mathfrak p_2\mathfrak p_3=(s).
\tag{A.185.15}
\]
The first prime is not principal, since \(a^2+6b^2=2\) has
no integral solution.  These relations prove that the class
group has order two and both prime classes generate it.
Inverting these primes therefore makes the class group trivial.
The principal divisor lattice supported there is exactly
\(\{(m,n)\in\Z^2:m+n\text{ is even}\}\); its basis is
\((2,0),(1,1)\), represented by \(2,s\).  Since the units
are \(\{\pm1\}\), the candidate squareclasses are
\((-1)^\sigma2^es^f\).  More explicitly, any class with
even valuations away from the two primes can first be divided
by a square so that its divisor is supported there: triviality
of the localized class group supplies the element whose
outside valuations are half of the original valuations.

The scaled dual to \(E_{8s}\) is \(E_{-2s}\), using
\((x,y)\mapsto(4x,8y)\) to \(E_{-32s}\).  At
\(\mathfrak p_3\), both coefficients have valuation one,
and the same two residue tests in \(\F_3\) give
\[
\begin{aligned}
 \alpha_{8s}(E_{8s}(k))&\subset\{1,-2,-s,2s\},\\
 \alpha_{-2s}(E_{-2s}(k))&\subset\{1,-2,s,-2s\}.
\end{aligned}
\tag{A.185.16}
\]
The point \((-2,-4+2s)\) gives \(-2\) in the first image,
and its kernel point gives \(2s\).  On the second curve
the actual point
\[
        (-5+4s,37-s),\qquad
        -5+4s=-2(2-s/2)^2
\tag{A.185.17}
\]
gives \(-2\), while the kernel point gives \(-2s\).
The class \(-2\) is not a square in \(k\): expanding
\((a+bs)^2=-2\) forces \(ab=0\), and then either
\(a^2=-2\) or \(b^2=1/3\), both impossible over \(\Q\).
The valuation at \(\mathfrak p_3\) separates its class
from the classes involving \(s\).  Both images have order
four, proving rank two.

\emph{The three original rational elliptics.}
Over \(\Q\), the candidate classes for all six coefficients
in the next table are \(\pm1,\pm2,\pm3,\pm6\).
Real signs and the primitive quartic (A.185.11) give exactly
the following upper bounds.  The third column records the
complete additional dyadic exclusions after the real ones.
\[
\begin{array}{r|l|l}
 b&\text{surviving classes}&\text{excluded classes and modulus}\\ \hline
 -6&1,3,-2,-6&-3,-1,2,6\pmod8\\
 24&1,6&2,3\pmod{16}\\
 18&1,2,3,6&\text{none}\\
 -72&1,-2&-6,-3,-1,2,3\pmod{16};\ 6\pmod{32}\\
 -12&1,-3,6,-2&-6,-1,2,3\pmod{16}\\
 48&1,3&2,6\pmod{16}.
\end{array}
\tag{A.185.18}
\]
In the positive-coefficient rows every negative class is
excluded over \(\R\).  Each finite exclusion means that,
for every residue pair \((U,V)\) with at least one odd
coordinate, \(d_0U^4+(b/d_0)V^4\) is not a square modulo
the displayed modulus.  This is a complete finite check of
at most \(32^2\) pairs, retained in the certificate.
The actual points \((3,3),(6,18),(6,12)\), together with
\((0,0)\), attain the four classes in the original rows.
The kernel point attains the nontrivial class in every dual
row.  Thus every upper bound is the exact connecting image,
and (A.185.12) gives rank one in all three cases.

Finally, (A.185.6), the actual scalings, and
\((\operatorname{Res}_{L/\Q}E)(\Q)=E(L)\) identify the
three Prym ranks with the computed elliptic ranks.  An isogeny
has an inverse homomorphism up to multiplication by its degree;
after tensoring the finitely generated rational point groups
with \(\Q\), it is an isomorphism.  Applying this to the
actual isogeny (A.182.3) proves (A.185.2).

The local, reproducible certificate
\path{code/point3_fermat8_prym_rank_certificate.py} verifies
all quotient formulas, differential scalars, field scalings,
point identities, and exact squareclass lists.  It also checks
the real-quadratic quartics modulo \((\sqrt2)^{13}\) and
reruns \path{code/point3_fermat8_elliptic_rank_certificate.py}.
The dyadic real-quadratic check is supplementary: the real
and odd-prime tests in the proof already give its complete
upper bounds, and displayed global points attain every class.
Retained output and JSON data specify all finite tests.
\end{proof}
